\section{Classical Mergeable-Sketch Parity: HyperLogLog}\label{sec:hll-parity}

The Markov walk-through showed that a tree of FNO leaves can handle an
oracle that defeats unordered register-wise merges without boundary
state. This section closes the loop in the opposite direction: when the oracle \emph{is} a
classical sketch's query, does the C-TreePO framework recover the
classical sketch exactly? The answer is yes --- up to byte-identity for
a register-based implementation --- and the register-space constraint is
itself informative about what learned operators can and cannot do.

We run the experiment on HyperLogLog \citep{FlajoletEtAl2007,HeuleEtAl2013}
because it is the canonical sketch for distinct counting and because the
merge operator (register-wise max) is associative at the state level. It is
also commutative and idempotent, but those are HLL-specific properties rather
than generic requirements of Proposition~\ref{prop:mergeable-reduction}.
HLL should be read as a representative member of the larger mergeable-sketch
class, not as the whole classical baseline: theorem-backed witnesses for
Count-Min, KLL/GK quantile, bigram, and Markov-count sketches are available
where their local-law obligations hold, and the crosswalk in
Appendix~\ref{app:proof-artifacts} lists each.

\subsection{Protocol}

All runs go through the unified \texttt{fit()} entry point. A preset
\texttt{classical\_hll\_parity\_task} builds a \texttt{TrainerConfig} whose
trainer is a zero-optimization callable that plugs a \texttt{SketchAdapter}
into \texttt{treepo\_reduce}. The same CSV schema holds four methods:

\begin{itemize}[leftmargin=1.5em]
    \item \textbf{classical-native:} the repo's own register-based HLL
          implementation (\texttt{HyperLogLogSketch}). Merge is pointwise
          max on a \texttt{uint8} register array.
    \item \textbf{classical-DataSketches:} Apache DataSketches'
          \texttt{hll\_sketch}/\texttt{hll\_union}, the canonical streaming
          reference. Internal state transitions through list, sparse, and
          dense modes based on cardinality.
    \item \textbf{learned-$g$ (classical $f$):} a learned neural merge
          $g$ produces a vector of dimension $m = 2^p$ interpreted as an
          HLL register state; the readout $f$ is the classical HLL
          estimator formula $\alpha_m m^2 / \sum 2^{-r_j}$ applied
          differentiably to that state.
    \item \textbf{learned-$g$+$f$ (end-to-end):} both the merge $g$ and
          the scalar readout $f$ are learned MLPs, with
          $\text{state\_dim} = 2\,\text{summary\_dim}$ respecting the
          FNO head 2$\times$ invariant documented in
          Appendix~\ref{app:mechanism-checks}.
\end{itemize}

For each method we sweep precisions $p \in \{7, 9, 11, 13\}$, leaf counts
$L \in \{1, 2, 4, 8, 16\}$, and three seeds. Learned runs train for $150$
epochs against per-node classical-HLL targets (\texttt{oracle\_kind} is
set to the classical HLL's own scoring head so the C1 (leaf sufficiency)
and C3 (merge consistency) targets equal what the flat HLL would produce
at each node; classical HLL's register-wise structure makes C2
(idempotence) automatic). In the projection language of
Section~\ref{sec:theorems}, classical HLL \emph{already} lies in the
exact local-law subspace by construction, so the learned runs are
projecting a learned-operator class onto the same subspace the
hand-designed sketch inhabits. Full reproduction details,
training-budget analysis, and the complete per-cell table are in
Appendix~\ref{app:classical-parity}; we summarize the headline findings
here.

\subsection{Proposition~\ref{prop:mergeable-reduction} Empirically}

For the classical methods, Proposition~\ref{prop:mergeable-reduction}
predicts that the tree-reduced summary should match the flat-reduced
summary up to oracle equivalence. Table~\ref{tab:classical-parity-hll}
confirms this on every cell of the grid:

\begin{itemize}[leftmargin=1.5em]
    \item \textbf{Native HLL: byte-exact.} The register vectors produced
          by $\texttt{reduce\_hll\_sketches}$ at every tree shape are
          byte-identical to those produced by the flat pipeline. The
          \texttt{bytes=} column is $1.00$ on every row.
    \item \textbf{DataSketches HLL: estimate-exact.} The \texttt{state\_equal}
          check (estimate agreement within $2\times$ HLL RSE) holds on
          every cell; the \texttt{bytes=} check fails because
          DataSketches' list$\to$sparse$\to$dense transitions depend on
          build order, but the functional oracle answer agrees.
\end{itemize}

This is Proposition~\ref{prop:mergeable-reduction}'s state-level collapse
made concrete. The native run instantiates $g$ as serialized HLL register
state and the merge as register-wise max; for a register-based implementation
the reduction is literally byte-for-byte.

\subsection{Can a Learned Merge Recover the Classical Sketch?}

The two learned methods ask a different question: if we pointed the same
\texttt{fit()} pipeline at the classical HLL as an oracle, how close can
a learned $g$ --- with or without a learned readout --- come to
reproducing its behavior? In the projection view, the classical methods
begin in the exact local-law subspace, while the learned methods test how
well different representable subspaces and readouts can approach that same
target under local-law supervision. Figure~\ref{fig:classical-parity-hll} shows
the relative MAE of each method against precision, one panel per leaf
count, and the picture tells a clean story:

\begin{itemize}[leftmargin=1.5em]
    \item \textbf{Learned $g$+$f$ (end-to-end)} tracks the classical
          sketch within a small multiple of the theoretical RSE floor and
          \emph{improves} as the tree gets deeper. At $p=13$, relative
          MAE drops from $\sim\!0.27$ at $L=2$ to $\sim\!0.19$ at
          $L \ge 8$: more merges mean more gradient signal at internal
          nodes, and the learned readout can compensate for imprecision
          anywhere in the state.
    \item \textbf{Learned $g$ with classical $f$} \emph{degrades} as the
          tree gets deeper. At $p=13$, relative MAE grows from $\sim\!0.27$
          at $L=1$ to $\sim\!0.43$ at $L=2$ to $\sim\!0.86$ at $L=16$.
          When the readout is a fixed non-learnable formula, mismatches
          at internal nodes channel through an unforgiving estimator and
          compound rather than average.
\end{itemize}

\begin{figure}[t]
    \centering
    \includegraphics[width=0.95\linewidth]{hll_parity_curves.png}
    \caption{Relative MAE against the oracle across precision, overlaying four methods on HyperLogLog: classical-native, classical-DataSketches, learned-$g$ (with classical $f$), and learned-$g$+$f$ (end-to-end). Dotted line: the HLL theoretical relative standard error $1.04/\sqrt{2^p}$. Each panel is a distinct leaf count $L \in \{1, 2, 4, 8, 16\}$. The two classical curves track each other within a fraction of the RSE floor, confirming Proposition~\ref{prop:mergeable-reduction}'s collapse. The learned-$g$+$f$ curve improves with $L$; the learned-$g$ (classical $f$) curve degrades with $L$.}\label{fig:classical-parity-hll}
\end{figure}

\begin{table}[t]
    \centering
    \small
    % Auto-generated by unified_g_v1.sketch.classical_parity_report; do not edit.
\begin{tabular}{llrrrrrrrrrrrrr}
\toprule
method & oracle & $p$ & $L$ & RSE & $\overline{|\Delta_{rel}|}$ & 95\% CI & $\max|\Delta|$ & state$=$ & bytes$=$ & rel MAE & C1 & C3 & wall$\times$ & $n$ \\
\midrule
classical\_datasketches & analytic & 11 & 1 & 0.02298 & 0 & 0 & 0 & 1 & 1 & 0.01133 & 15.14 & 0 & 0.8513 & 3 \\
classical\_datasketches & analytic & 11 & 16 & 0.02298 & 0.008052 & 0.0008999 & 13.05 & 1 & 0 & 0.01322 & 0.07798 & 8.668 & 1.832 & 3 \\
classical\_datasketches & analytic & 11 & 2 & 0.02298 & 0.008025 & 0.0008882 & 13.05 & 1 & 0 & 0.01319 & 7.093 & 17.24 & 1.134 & 3 \\
classical\_datasketches & analytic & 11 & 4 & 0.02298 & 0.008054 & 0.000927 & 13.07 & 1 & 0 & 0.01322 & 2.854 & 12.88 & 1.319 & 3 \\
classical\_datasketches & analytic & 11 & 8 & 0.02298 & 0.008053 & 0.0009013 & 13.05 & 1 & 0 & 0.01317 & 0.885 & 10.13 & 1.577 & 3 \\
classical\_datasketches & analytic & 13 & 1 & 0.01149 & 0 & 0 & 0 & 1 & 1 & 0.002884 & 5.227 & 0 & 1.032 & 3 \\
classical\_datasketches & analytic & 13 & 16 & 0.01149 & 0.002351 & 0.0008971 & 5.651 & 1 & 0 & 0.003686 & 3.394e-05 & 1.994 & 1.572 & 3 \\
classical\_datasketches & analytic & 13 & 2 & 0.01149 & 0.002365 & 0.0008783 & 5.53 & 1 & 0 & 0.003722 & 1.378 & 7.073 & 1.221 & 3 \\
classical\_datasketches & analytic & 13 & 4 & 0.01149 & 0.002415 & 0.0007943 & 5.538 & 1 & 0 & 0.003728 & 0.1452 & 3.593 & 1.376 & 3 \\
classical\_datasketches & analytic & 13 & 8 & 0.01149 & 0.002327 & 0.000767 & 5.452 & 1 & 0 & 0.003675 & 0.0009903 & 2.417 & 1.531 & 3 \\
classical\_datasketches & analytic & 7 & 1 & 0.09192 & 0 & 0 & 0 & 1 & 1 & 0.05798 & 73.93 & 0 & 1.053 & 3 \\
classical\_datasketches & analytic & 7 & 16 & 0.09192 & 0.04028 & 0.004683 & 59.86 & 1 & 0 & 0.07586 & 4.539 & 54.61 & 1.321 & 3 \\
classical\_datasketches & analytic & 7 & 2 & 0.09192 & 0.04028 & 0.004683 & 59.86 & 1 & 0 & 0.07586 & 35.33 & 102.8 & 1.073 & 3 \\
classical\_datasketches & analytic & 7 & 4 & 0.09192 & 0.04028 & 0.004683 & 59.86 & 1 & 0 & 0.07586 & 18.97 & 73.53 & 1.113 & 3 \\
classical\_datasketches & analytic & 7 & 8 & 0.09192 & 0.04028 & 0.004683 & 59.86 & 1 & 0 & 0.07586 & 9.047 & 61.06 & 0.9858 & 3 \\
classical\_datasketches & analytic & 9 & 1 & 0.04596 & 0 & 0 & 0 & 1 & 1 & 0.02289 & 29 & 0 & 1.115 & 3 \\
classical\_datasketches & analytic & 9 & 16 & 0.04596 & 0.01965 & 0.001969 & 31.12 & 1 & 0 & 0.02862 & 1.82 & 22.09 & 1.825 & 3 \\
classical\_datasketches & analytic & 9 & 2 & 0.04596 & 0.01965 & 0.001969 & 31.12 & 1 & 0 & 0.02862 & 15.27 & 39.39 & 1.085 & 3 \\
classical\_datasketches & analytic & 9 & 4 & 0.04596 & 0.01965 & 0.001969 & 31.12 & 1 & 0 & 0.02862 & 7.661 & 29.39 & 1.177 & 3 \\
classical\_datasketches & analytic & 9 & 8 & 0.04596 & 0.01965 & 0.001969 & 31.12 & 1 & 0 & 0.02862 & 3.956 & 24.61 & 1.321 & 3 \\
classical\_datasketches & hll\_reference & 11 & 1 & 0.02298 & 0 & 0 & 0 & 1 & 1 & 0 & 0 & 0 & 1.041 & 3 \\
classical\_datasketches & hll\_reference & 11 & 16 & 0.02298 & 0.008052 & 0.0008999 & 13.05 & 1 & 0 & 0.008052 & 0 & 2.508 & 1.83 & 3 \\
classical\_datasketches & hll\_reference & 11 & 2 & 0.02298 & 0.008025 & 0.0008882 & 13.05 & 1 & 0 & 0.008025 & 0 & 10.68 & 1.141 & 3 \\
classical\_datasketches & hll\_reference & 11 & 4 & 0.02298 & 0.008054 & 0.000927 & 13.07 & 1 & 0 & 0.008054 & 0 & 7.789 & 1.322 & 3 \\
classical\_datasketches & hll\_reference & 11 & 8 & 0.02298 & 0.008053 & 0.0009013 & 13.05 & 1 & 0 & 0.008053 & 0 & 5.765 & 1.513 & 3 \\
classical\_datasketches & hll\_reference & 13 & 1 & 0.01149 & 0 & 0 & 0 & 1 & 1 & 0 & 0 & 0 & 1.047 & 3 \\
classical\_datasketches & hll\_reference & 13 & 16 & 0.01149 & 0.002351 & 0.0008971 & 5.651 & 1 & 0 & 0.002351 & 0 & 0.6059 & 1.678 & 3 \\
classical\_datasketches & hll\_reference & 13 & 2 & 0.01149 & 0.002365 & 0.0008783 & 5.53 & 1 & 0 & 0.002365 & 0 & 4.251 & 1.186 & 3 \\
classical\_datasketches & hll\_reference & 13 & 4 & 0.01149 & 0.002415 & 0.0007943 & 5.538 & 1 & 0 & 0.002415 & 0 & 1.858 & 1.379 & 3 \\
classical\_datasketches & hll\_reference & 13 & 8 & 0.01149 & 0.002327 & 0.000767 & 5.452 & 1 & 0 & 0.002327 & 0 & 0.9663 & 1.535 & 3 \\
classical\_datasketches & hll\_reference & 7 & 1 & 0.09192 & 0 & 0 & 0 & 1 & 1 & 0 & 0 & 0 & 1.054 & 3 \\
classical\_datasketches & hll\_reference & 7 & 16 & 0.09192 & 0.04028 & 0.004683 & 59.86 & 1 & 0 & 0.04028 & 0 & 32.12 & 1.308 & 3 \\
classical\_datasketches & hll\_reference & 7 & 2 & 0.09192 & 0.04028 & 0.004683 & 59.86 & 1 & 0 & 0.04028 & 0 & 53.52 & 1.078 & 3 \\
classical\_datasketches & hll\_reference & 7 & 4 & 0.09192 & 0.04028 & 0.004683 & 59.86 & 1 & 0 & 0.04028 & 0 & 41.47 & 1.121 & 3 \\
classical\_datasketches & hll\_reference & 7 & 8 & 0.09192 & 0.04028 & 0.004683 & 59.86 & 1 & 0 & 0.04028 & 0 & 35.49 & 1.102 & 3 \\
classical\_datasketches & hll\_reference & 9 & 1 & 0.04596 & 0 & 0 & 0 & 1 & 1 & 0 & 0 & 0 & 1.052 & 3 \\
classical\_datasketches & hll\_reference & 9 & 16 & 0.04596 & 0.01965 & 0.001969 & 31.12 & 1 & 0 & 0.01965 & 0 & 14.5 & 1.568 & 3 \\
classical\_datasketches & hll\_reference & 9 & 2 & 0.04596 & 0.01965 & 0.001969 & 31.12 & 1 & 0 & 0.01965 & 0 & 27.06 & 0.8283 & 3 \\
classical\_datasketches & hll\_reference & 9 & 4 & 0.04596 & 0.01965 & 0.001969 & 31.12 & 1 & 0 & 0.01965 & 0 & 19.52 & 1.166 & 3 \\
classical\_datasketches & hll\_reference & 9 & 8 & 0.04596 & 0.01965 & 0.001969 & 31.12 & 1 & 0 & 0.01965 & 0 & 16.1 & 1.314 & 3 \\
classical\_native & analytic & 11 & 1 & 0.02298 & 0 & 0 & 0 & 1 & 1 & 0.01461 & 19.47 & 0 & 1.018 & 3 \\
classical\_native & analytic & 11 & 16 & 0.02298 & 0 & 0 & 0 & 1 & 1 & 0.01461 & 1.441 & 10.92 & 1.135 & 3 \\
classical\_native & analytic & 11 & 2 & 0.02298 & 0 & 0 & 0 & 1 & 1 & 0.01461 & 10.44 & 19.47 & 0.9988 & 3 \\
classical\_native & analytic & 11 & 4 & 0.02298 & 0 & 0 & 0 & 1 & 1 & 0.01461 & 5.081 & 14.54 & 1.041 & 3 \\
classical\_native & analytic & 11 & 8 & 0.02298 & 0 & 0 & 0 & 1 & 1 & 0.01461 & 2.753 & 12.23 & 1.061 & 3 \\
classical\_native & analytic & 13 & 1 & 0.01149 & 0 & 0 & 0 & 1 & 1 & 0.008445 & 9.396 & 0 & 0.9912 & 3 \\
classical\_native & analytic & 13 & 16 & 0.01149 & 0 & 0 & 0 & 1 & 1 & 0.008445 & 0.7089 & 5.498 & 1.094 & 3 \\
classical\_native & analytic & 13 & 2 & 0.01149 & 0 & 0 & 0 & 1 & 1 & 0.008445 & 5.149 & 9.396 & 1.025 & 3 \\
classical\_native & analytic & 13 & 4 & 0.01149 & 0 & 0 & 0 & 1 & 1 & 0.008445 & 2.723 & 7.212 & 1.017 & 3 \\
classical\_native & analytic & 13 & 8 & 0.01149 & 0 & 0 & 0 & 1 & 1 & 0.008445 & 1.405 & 6.127 & 1.074 & 3 \\
classical\_native & analytic & 7 & 1 & 0.09192 & 0 & 0 & 0 & 1 & 1 & 0.06938 & 84.16 & 0 & 0.9672 & 3 \\
classical\_native & analytic & 7 & 16 & 0.09192 & 0 & 0 & 0 & 1 & 1 & 0.06938 & 6.93 & 51 & 1.139 & 3 \\
classical\_native & analytic & 7 & 2 & 0.09192 & 0 & 0 & 0 & 1 & 1 & 0.06938 & 49.8 & 84.16 & 1.02 & 3 \\
classical\_native & analytic & 7 & 4 & 0.09192 & 0 & 0 & 0 & 1 & 1 & 0.06938 & 24.07 & 65.44 & 1.054 & 3 \\
classical\_native & analytic & 7 & 8 & 0.09192 & 0 & 0 & 0 & 1 & 1 & 0.06938 & 13.36 & 56.09 & 1.094 & 3 \\
classical\_native & analytic & 9 & 1 & 0.04596 & 0 & 0 & 0 & 1 & 1 & 0.03318 & 41.96 & 0 & 0.9943 & 3 \\
classical\_native & analytic & 9 & 16 & 0.04596 & 0 & 0 & 0 & 1 & 1 & 0.03318 & 2.744 & 24.37 & 1.12 & 3 \\
classical\_native & analytic & 9 & 2 & 0.04596 & 0 & 0 & 0 & 1 & 1 & 0.03318 & 23.93 & 41.96 & 1.05 & 3 \\
classical\_native & analytic & 9 & 4 & 0.04596 & 0 & 0 & 0 & 1 & 1 & 0.03318 & 11.5 & 31.98 & 1.07 & 3 \\
classical\_native & analytic & 9 & 8 & 0.04596 & 0 & 0 & 0 & 1 & 1 & 0.03318 & 5.598 & 27.01 & 1.069 & 3 \\
classical\_native & hll\_reference & 11 & 1 & 0.02298 & 0 & 0 & 0 & 1 & 1 & 0 & 0 & 0 & 1.009 & 3 \\
classical\_native & hll\_reference & 11 & 16 & 0.02298 & 0 & 0 & 0 & 1 & 1 & 0 & 0 & 0 & 1.096 & 3 \\
classical\_native & hll\_reference & 11 & 2 & 0.02298 & 0 & 0 & 0 & 1 & 1 & 0 & 0 & 0 & 1.017 & 3 \\
classical\_native & hll\_reference & 11 & 4 & 0.02298 & 0 & 0 & 0 & 1 & 1 & 0 & 0 & 0 & 1.066 & 3 \\
classical\_native & hll\_reference & 11 & 8 & 0.02298 & 0 & 0 & 0 & 1 & 1 & 0 & 0 & 0 & 1.03 & 3 \\
classical\_native & hll\_reference & 13 & 1 & 0.01149 & 0 & 0 & 0 & 1 & 1 & 0 & 0 & 0 & 0.9949 & 3 \\
classical\_native & hll\_reference & 13 & 16 & 0.01149 & 0 & 0 & 0 & 1 & 1 & 0 & 0 & 0 & 1.08 & 3 \\
classical\_native & hll\_reference & 13 & 2 & 0.01149 & 0 & 0 & 0 & 1 & 1 & 0 & 0 & 0 & 1.007 & 3 \\
classical\_native & hll\_reference & 13 & 4 & 0.01149 & 0 & 0 & 0 & 1 & 1 & 0 & 0 & 0 & 1.033 & 3 \\
classical\_native & hll\_reference & 13 & 8 & 0.01149 & 0 & 0 & 0 & 1 & 1 & 0 & 0 & 0 & 1.011 & 3 \\
classical\_native & hll\_reference & 7 & 1 & 0.09192 & 0 & 0 & 0 & 1 & 1 & 0 & 0 & 0 & 0.9832 & 3 \\
classical\_native & hll\_reference & 7 & 16 & 0.09192 & 0 & 0 & 0 & 1 & 1 & 0 & 0 & 0 & 1.158 & 3 \\
classical\_native & hll\_reference & 7 & 2 & 0.09192 & 0 & 0 & 0 & 1 & 1 & 0 & 0 & 0 & 1.05 & 3 \\
classical\_native & hll\_reference & 7 & 4 & 0.09192 & 0 & 0 & 0 & 1 & 1 & 0 & 0 & 0 & 1.069 & 3 \\
classical\_native & hll\_reference & 7 & 8 & 0.09192 & 0 & 0 & 0 & 1 & 1 & 0 & 0 & 0 & 1.11 & 3 \\
classical\_native & hll\_reference & 9 & 1 & 0.04596 & 0 & 0 & 0 & 1 & 1 & 0 & 0 & 0 & 0.9869 & 3 \\
classical\_native & hll\_reference & 9 & 16 & 0.04596 & 0 & 0 & 0 & 1 & 1 & 0 & 0 & 0 & 1.146 & 3 \\
classical\_native & hll\_reference & 9 & 2 & 0.04596 & 0 & 0 & 0 & 1 & 1 & 0 & 0 & 0 & 1.052 & 3 \\
classical\_native & hll\_reference & 9 & 4 & 0.04596 & 0 & 0 & 0 & 1 & 1 & 0 & 0 & 0 & 1.032 & 3 \\
classical\_native & hll\_reference & 9 & 8 & 0.04596 & 0 & 0 & 0 & 1 & 1 & 0 & 0 & 0 & 1.076 & 3 \\
learned\_g & hll\_reference & 11 & 1 & 0.02298 & -- & -- & -- & -- & -- & 0.3265 & 66.2 & 0 & -- & 3 \\
learned\_g & hll\_reference & 11 & 16 & 0.02298 & -- & -- & -- & -- & -- & 0.8796 & 3.599 & 95.57 & -- & 3 \\
learned\_g & hll\_reference & 11 & 2 & 0.02298 & -- & -- & -- & -- & -- & 0.5141 & 31.86 & 105.6 & -- & 3 \\
learned\_g & hll\_reference & 11 & 4 & 0.02298 & -- & -- & -- & -- & -- & 0.4717 & 13.91 & 70.56 & -- & 3 \\
learned\_g & hll\_reference & 11 & 8 & 0.02298 & -- & -- & -- & -- & -- & 0.7468 & 6.413 & 87.07 & -- & 3 \\
learned\_g & hll\_reference & 13 & 1 & 0.01149 & -- & -- & -- & -- & -- & 0.2866 & 58.25 & 0 & -- & 3 \\
learned\_g & hll\_reference & 13 & 16 & 0.01149 & -- & -- & -- & -- & -- & 0.8893 & 3.847 & 98.03 & -- & 3 \\
learned\_g & hll\_reference & 13 & 2 & 0.01149 & -- & -- & -- & -- & -- & 0.4998 & 35.13 & 103.7 & -- & 3 \\
learned\_g & hll\_reference & 13 & 4 & 0.01149 & -- & -- & -- & -- & -- & 0.7436 & 16.13 & 107.5 & -- & 3 \\
learned\_g & hll\_reference & 13 & 8 & 0.01149 & -- & -- & -- & -- & -- & 0.7201 & 7.373 & 82.79 & -- & 3 \\
learned\_g & hll\_reference & 7 & 1 & 0.09192 & -- & -- & -- & -- & -- & 0.3292 & 65.61 & 0 & -- & 3 \\
learned\_g & hll\_reference & 7 & 16 & 0.09192 & -- & -- & -- & -- & -- & 0.186 & 8.146 & 27.95 & -- & 3 \\
learned\_g & hll\_reference & 7 & 2 & 0.09192 & -- & -- & -- & -- & -- & 0.2623 & 29.34 & 52.37 & -- & 3 \\
learned\_g & hll\_reference & 7 & 4 & 0.09192 & -- & -- & -- & -- & -- & 0.1863 & 14.18 & 33.83 & -- & 3 \\
learned\_g & hll\_reference & 7 & 8 & 0.09192 & -- & -- & -- & -- & -- & 0.1778 & 8.486 & 28.07 & -- & 3 \\
learned\_g & hll\_reference & 9 & 1 & 0.04596 & -- & -- & -- & -- & -- & 0.2654 & 53.36 & 0 & -- & 3 \\
learned\_g & hll\_reference & 9 & 16 & 0.04596 & -- & -- & -- & -- & -- & 0.4396 & 4.974 & 53.75 & -- & 3 \\
learned\_g & hll\_reference & 9 & 2 & 0.04596 & -- & -- & -- & -- & -- & 0.2657 & 30.51 & 53.77 & -- & 3 \\
learned\_g & hll\_reference & 9 & 4 & 0.04596 & -- & -- & -- & -- & -- & 0.2919 & 16.34 & 44.86 & -- & 3 \\
learned\_g & hll\_reference & 9 & 8 & 0.04596 & -- & -- & -- & -- & -- & 0.2835 & 8.893 & 39.25 & -- & 3 \\
learned\_fg & hll\_reference & 11 & 1 & 0.02298 & -- & -- & -- & -- & -- & 0.2963 & 60.48 & 0 & -- & 3 \\
learned\_fg & hll\_reference & 11 & 16 & 0.02298 & -- & -- & -- & -- & -- & 0.2481 & 5.789 & 31.28 & -- & 3 \\
learned\_fg & hll\_reference & 11 & 2 & 0.02298 & -- & -- & -- & -- & -- & 0.2488 & 29.12 & 50.53 & -- & 3 \\
learned\_fg & hll\_reference & 11 & 4 & 0.02298 & -- & -- & -- & -- & -- & 0.2465 & 14.89 & 39.7 & -- & 3 \\
learned\_fg & hll\_reference & 11 & 8 & 0.02298 & -- & -- & -- & -- & -- & 0.223 & 8.897 & 33.8 & -- & 3 \\
learned\_fg & hll\_reference & 13 & 1 & 0.01149 & -- & -- & -- & -- & -- & 0.2649 & 53.93 & 0 & -- & 3 \\
learned\_fg & hll\_reference & 13 & 16 & 0.01149 & -- & -- & -- & -- & -- & 0.2489 & 7.425 & 31.66 & -- & 3 \\
learned\_fg & hll\_reference & 13 & 2 & 0.01149 & -- & -- & -- & -- & -- & 0.2657 & 30.16 & 54.39 & -- & 3 \\
learned\_fg & hll\_reference & 13 & 4 & 0.01149 & -- & -- & -- & -- & -- & 0.2366 & 13.55 & 36.4 & -- & 3 \\
learned\_fg & hll\_reference & 13 & 8 & 0.01149 & -- & -- & -- & -- & -- & 0.1825 & 7.409 & 34.55 & -- & 3 \\
learned\_fg & hll\_reference & 7 & 1 & 0.09192 & -- & -- & -- & -- & -- & 0.298 & 59.83 & 0 & -- & 3 \\
learned\_fg & hll\_reference & 7 & 16 & 0.09192 & -- & -- & -- & -- & -- & 0.2491 & 6.522 & 34.85 & -- & 3 \\
learned\_fg & hll\_reference & 7 & 2 & 0.09192 & -- & -- & -- & -- & -- & 0.2726 & 29.94 & 53.97 & -- & 3 \\
learned\_fg & hll\_reference & 7 & 4 & 0.09192 & -- & -- & -- & -- & -- & 0.2514 & 13.81 & 38 & -- & 3 \\
learned\_fg & hll\_reference & 7 & 8 & 0.09192 & -- & -- & -- & -- & -- & 0.2326 & 7.277 & 37.82 & -- & 3 \\
learned\_fg & hll\_reference & 9 & 1 & 0.04596 & -- & -- & -- & -- & -- & 0.2949 & 59.31 & 0 & -- & 3 \\
learned\_fg & hll\_reference & 9 & 16 & 0.04596 & -- & -- & -- & -- & -- & 0.2485 & 6.742 & 36.4 & -- & 3 \\
learned\_fg & hll\_reference & 9 & 2 & 0.04596 & -- & -- & -- & -- & -- & 0.2687 & 30.47 & 54.14 & -- & 3 \\
learned\_fg & hll\_reference & 9 & 4 & 0.04596 & -- & -- & -- & -- & -- & 0.2599 & 14.64 & 40.82 & -- & 3 \\
learned\_fg & hll\_reference & 9 & 8 & 0.04596 & -- & -- & -- & -- & -- & 0.2054 & 10.76 & 41.49 & -- & 3 \\
\bottomrule
\end{tabular}

    \caption{HLL parity summary across the \texttt{(method, oracle, precision, $L$)} grid. \texttt{RSE} is the HLL theoretical relative standard error $1.04/\sqrt{2^p}$. \texttt{state=} and \texttt{bytes=} are per-cell equality rates for the classical rows; \texttt{rel MAE} is each method's relative mean absolute error against the oracle; \texttt{C1} and \texttt{C3} are the per-leaf and per-merge MAE; \texttt{wall$\times$} is the tree-to-flat wall-time ratio. Full table with all seeds and oracle-kind variants in Appendix~\ref{app:classical-parity}.}\label{tab:classical-parity-hll}
\end{table}

\subsection{Takeaway}

Two structural findings follow from this experiment:

\paragraph{The classical collapse is exact.}
When the state-level conditions of Proposition~\ref{prop:mergeable-reduction} hold,
C-TreePO reduces to the classical mergeable sketch not merely in
distribution but byte-for-byte on the reference implementation. This is
the tightest empirical version of the generalization claim the paper can
make for this representative sketch; the formal version is the
sketch-generic inclusion
\texttt{mergeableSketchSummaryClass\_subset\_exactLocalLawSubspace}.

\paragraph{The register-space constraint is a real optimization barrier.}
With identical data, identical training objective, and identical
hyperparameters, the only difference between \texttt{learned\_g} and
\texttt{learned\_gf} is whether the readout $f$ is learnable or fixed to
the classical HLL estimator. The gap between the two --- and its
direction-of-scaling with $L$ --- is therefore a clean statement that
some learned-$g$ settings pay a real cost when the readout must live in
the classical sketch's exact state space. A practitioner building on this
framework who does not need byte-for-byte parity with a specific
reference sketch would pick the end-to-end learned path.

The Markov interlude established that the framework handles oracles no
classical sketch can, and this section shows it recovers the classical
mergeable-sketch route when one exists. Section~\ref{sec:theorems} now states the
theorems that govern both: local validity composing to root
preservation, preservation lifting to training-equivalence for DPO and
GRPO, the oracle/projection view that turns local-law weights into a
structured approximation-error assumption, and sampled audits yielding
finite-sample gap bounds. Once that
machinery is in place, Section~\ref{sec:manifesto-llm} applies it to a
setting where no classical sketch exists --- the Manifesto-Project RILE
rubric --- and the learned $g$ must carry the full semantic load.
